\begin{center}
\LARGE
Application of Groebner Bases for Numerical Constraint
                             Solving

\end{center}

\begin{center}
\Large
Frederic Benhamou, Laurent Granvilliers

\end{center}

\noindent
Date:  July 19th (Friday)\\
Time:  11:00-11:30\\

\begin{center}
\large
Abstract
\end{center}

This work is an attempt to improve the computational behaviour of interval propagation methods
by using techniques from computer algebra such as Groebner bases for solving polynomial
constraints over the real numbers. 

Interval propagation methods implement local consistency techniques and have to be combined
together with enumeration techniques to compute approximations of the solutions satisfying the
initial constraint set. Unfortunately the number of enumeration steps may be important due to
unappropriate expressions of the initial problem. We show how Groebner bases computations,
used as a preprocessing step, may generate better expressions of the initial problem which can be
solved faster using interval propagation methods. Furthermore, we provide a general semantics for
combining these different methods. 

In the talk we present some small examples to illustrate the method and we give experimental
results. 

